\paragraph{}
The use of \gls{ml} in medical prediction creates two related challenges: published predictive workflows may be difficult to reproduce, and post-hoc explanations may vary according to the model, data representation, and explanation method. This dissertation developed and evaluated a reproducible \gls{xai} framework through the replication and extension of five medical-data studies concerning heart disease, stroke, diabetes, \gls{ckd}, and \gls{hbaone}-associated hospital readmission.

The source workflows were reconstructed in dataset-specific notebooks. Paper-style replication was distinguished from corrected analytical branches whenever preprocessing order, information leakage, class imbalance, or evaluation design limited the validity of direct imitation. Four final models were evaluated on held-out test data; the \gls{hbaone} study retained its original full-cohort inferential design. A common framework then generated native model importance where supported, permutation importance, \gls{shap}, repeated \gls{lime} explanations for up to ten observations per case category and 20 perturbation seeds, rank-agreement measures, and bootstrap-refit stability results when retraining was available.

Global-method agreement was strongly dataset- and representation-dependent. Native importance and \gls{shap} were closely aligned for Heart, Stroke, Diabetes, and \gls{ckd}, whereas permutation rankings diverged more often. In the interaction-rich \gls{hbaone} model, \gls{shap} and permutation importance agreed strongly, while absolute coefficient magnitude showed almost no agreement with \gls{shap}. Across 100 bootstrap refits, complete-ranking stability was highest for \gls{ckd}, followed by Diabetes and Stroke, while Heart showed lower but still positive persistence. Diabetes had the highest overlap among the three leading features. Mean seed-to-seed \gls{lime} rank correlations ranged from 0.898 to 0.969, yet correct and incorrect predictions frequently shared the same strongest local rules, including confident Stroke errors. Consequently, coherent and repeatable explanations did not independently establish prediction correctness.

The framework executed successfully across linear classifiers, tree ensembles, and the tested logistic-regression implementation. It was supported by dependency locks, provenance metadata, saved bundles, command-line reruns, comparison reports, and automated tests. Its generalisability remains bounded to the tested binary-classification interfaces. No prospective, causal, clinician-centred, or clinical-utility validation was performed. The principal conclusion is that \gls{xai} should be treated as a reproducible, multi-method model-audit process rather than as an automatic certificate of clinical trust.

\paragraph{}
\textbf{Keywords:} Explainable Artificial Intelligence; Clinical Prediction; Reproducibility; SHAP; LIME; Feature Importance; Explanation Stability.
